% Shared figure block, \input by BOTH main.tex (via results/frontier.tex) and
% notebook.tex (via experiments.tex). The two are independent builds, so one
% label here cannot clash; keeping a single copy means the caption cannot drift
% between the paper and the record, which is exactly what happened to the CLEO
% frontier row before the split.
\begin{figure}[htbp]
  \centering
  \includegraphics[width=\textwidth]{ame_distributions}
  \caption{\textbf{Where the distributions actually sit, rather than the success
  rate thresholded out of them.} Each row is one arm's per-design best-of-5 motif
  RMSD: box for median and quartiles, jittered strip behind it for the shape the
  box hides. Log axis, because a linear one lets the single worst arm compress
  everything that matters into the left fifth. \textbf{Left, M0097:} every
  \cleo{} arm has its \emph{median} below the cutoff (\SIrange{1.07}{1.44}{\angstrom})
  --- the mode has moved, not just the tail --- against
  \SIrange{1.31}{1.42}{\angstrom} for low-temperature LigandMPNN and
  \SI{5.59}{\angstrom} at $T{=}1.0$. The non-monotonicity noted in E10 is visible
  directly: $T{=}0.3$ has a \emph{lower} median than $T{=}0.1$.
  \textbf{Right, the three near-miss backbones:} medians of
  \SIrange{2.81}{3.12}{\angstrom} with the cutoff far out in the left tail. The
  one design that survived 40 refolds sits at the extreme edge of its own
  distribution. Encoding: both ladders are ordinal, so each is a single hue
  ramped light to dark rather than arbitrary categorical colours; the anchored
  arm keeps an accent hue as the highlighted result.}
  \label{fig:ame-distributions}
\end{figure}
